% Options for packages loaded elsewhere
\PassOptionsToPackage{unicode}{hyperref}
\PassOptionsToPackage{hyphens}{url}
\documentclass[
]{article}
\usepackage{xcolor}
\usepackage{amsmath,amssymb}
\setcounter{secnumdepth}{-\maxdimen} % remove section numbering
\usepackage{iftex}
\ifPDFTeX
  \usepackage[T1]{fontenc}
  \usepackage[utf8]{inputenc}
  \usepackage{textcomp} % provide euro and other symbols
\else % if luatex or xetex
  \usepackage{unicode-math} % this also loads fontspec
  \defaultfontfeatures{Scale=MatchLowercase}
  \defaultfontfeatures[\rmfamily]{Ligatures=TeX,Scale=1}
\fi
\usepackage{lmodern}
\ifPDFTeX\else
  % xetex/luatex font selection
\fi
% Use upquote if available, for straight quotes in verbatim environments
\IfFileExists{upquote.sty}{\usepackage{upquote}}{}
\IfFileExists{microtype.sty}{% use microtype if available
  \usepackage[]{microtype}
  \UseMicrotypeSet[protrusion]{basicmath} % disable protrusion for tt fonts
}{}
\makeatletter
\@ifundefined{KOMAClassName}{% if non-KOMA class
  \IfFileExists{parskip.sty}{%
    \usepackage{parskip}
  }{% else
    \setlength{\parindent}{0pt}
    \setlength{\parskip}{6pt plus 2pt minus 1pt}}
}{% if KOMA class
  \KOMAoptions{parskip=half}}
\makeatother
\usepackage{longtable,booktabs,array}
\newcounter{none} % for unnumbered tables
\usepackage{calc} % for calculating minipage widths
% Correct order of tables after \paragraph or \subparagraph
\usepackage{etoolbox}
\makeatletter
\patchcmd\longtable{\par}{\if@noskipsec\mbox{}\fi\par}{}{}
\makeatother
% Allow footnotes in longtable head/foot
\IfFileExists{footnotehyper.sty}{\usepackage{footnotehyper}}{\usepackage{footnote}}
\makesavenoteenv{longtable}
\setlength{\emergencystretch}{3em} % prevent overfull lines
\providecommand{\tightlist}{%
  \setlength{\itemsep}{0pt}\setlength{\parskip}{0pt}}
\usepackage{bookmark}
\IfFileExists{xurl.sty}{\usepackage{xurl}}{} % add URL line breaks if available
\urlstyle{same}
\hypersetup{
  hidelinks,
  pdfcreator={LaTeX via pandoc}}

\author{}
\date{}

\begin{document}

\section{Collateral Damage as Narrative Convention: False Accusation and
Its Costs Across 27 Seasons of Law \& Order: Special Victims
Unit}\label{collateral-damage-as-narrative-convention-false-accusation-and-its-costs-across-27-seasons-of-law-order-special-victims-unit}

\textbf{Tyler Satchel Orden} Independent researcher Contact:
tylerorden@gmail.com

\emph{Manuscript prepared July 2026}

\begin{center}\rule{0.5\linewidth}{0.5pt}\end{center}

\subsection{Abstract}\label{abstract}

Law \& Order: Special Victims Unit is one of the most durable fictions
about policing ever broadcast in the United States. This article reports
a census-style content analysis of the show's 27 seasons, asking how
often the series depicts innocent people harmed by false or mistaken
accusations, what those harms look like, and how the fictional
detectives answer for them. A large language model applied a fixed
codebook to every episode transcript, functioning as a measurement
instrument; a ground-truth audit then found that nine source cells held
another episode's script, a silent corruption the coding had inherited,
and all analyses use the corrected set of 567 episodes and 521
person-episode records. At least one falsely accused person appears in
59.4\% of episodes (62.1\% counting ambiguous cases), 62.4\% of persons
reach severity 3 or 4 on a four-point scale, and a row-level audit,
adjudicated case by case, finds 32 accusation-linked deaths: 21 murders,
9 suicides, and 2 others. The detectives' own inference is the leading
accusation origin (54.5\%), detectives threaten, coerce, or degrade the
innocent in 53.0\% of cases, and where an accusation spreads beyond the
precinct the squad itself is the discloser 83.9\% of the time. Yet an
apology of any kind follows in only 7.5\% of cases, and a formal apology
in 1.5\%. Apology tracks catastrophe rather than conduct: 0.0\% at
severity 1 against 16.4\% at severity 4, while threatened and
unthreatened persons are apologized to at nearly identical rates. The
show also pairs media exposure with catastrophe: 45.3\% of media-exposed
persons reach severity 4, against 8.5\% when word never leaves the
precinct. I argue that the series renders collateral harm to the
innocent as a narrative convention, visible but unaccounted for.

\textbf{Keywords:} false accusations; police procedural; copaganda;
cultivation theory; content analysis; Law \& Order: Special Victims
Unit; large language models; wrongful conviction

\begin{center}\rule{0.5\linewidth}{0.5pt}\end{center}

\subsection{1. Introduction}\label{introduction}

In the eighth episode of its second season, Law \& Order: Special
Victims Unit arrests a man named Russell Ramsay for a rape he did not
commit. The accusation reaches him through the sex-offender registry. He
is perp-walked, charged, and remanded to Rikers Island, where he begs
the detectives for help: ``I gotta stay awake at night so I don't get
ambushed\ldots{} I can't look at a guy without getting jumped.'' A few
scenes later the outcome arrives in three flat sentences: ``Russell
Ramsay was attacked at Rikers. He was beaten and raped. He was killed.''
The episode then identifies the real rapist and ends.

Ramsay is fictional, and so is the harm the script assigns him. What
interests me is that the series has done some version of this to 520
other people.

SVU premiered in September 1999 and was still airing new episodes when I
collected data in January 2026: 27 seasons, 576 episodes, plus a second
life in syndication and streaming that puts its reruns in front of
viewers more or less continuously. Scholarship on the show has
concentrated on its treatment of victims and its gender politics, and a
larger literature treats the police procedural as a machine for
producing favorable beliefs about policing; Section 2 reviews both.
Almost none of that work looks squarely at the people the fictional
investigation chews through on its way to the real perpetrator: the
suspects who are accused, exposed, and then cleared.

That omission is worth repairing, because the cleared suspect is where a
procedural quietly states its theory of acceptable loss. A show can be
sympathetic to victims, hostile to offenders, and still teach its
audience something specific about the innocent people in between: how
often they get caught in the machinery, how badly they are hurt, who
pointed the finger first, and whether anyone ever says sorry. Those are
empirical questions about the text itself, and they can be answered by
counting.

This study counts.

I coded every SVU episode broadcast through January 2026 for persons who
are accused of a crime, later shown or strongly implied to be innocent,
and harmed by the accusation through public exposure or physical
consequence. The coding was performed by a large language model applying
a fixed written codebook, a design I describe in Section 3 and treat
throughout as an instrument to be audited rather than a collaborator to
be trusted. Auditing mattered more than I planned: it caught corruption
in the source transcripts themselves, and every figure here comes from
the corrected set, 521 person-episode records across 567 verified
episodes. One result belongs up front: the row-level death audit finds
32 accusation-linked deaths, among them 21 murders and 9 suicides.

The pattern the full results describe, I will argue, is a narrative
convention: collateral damage to the innocent as a routine, visible, and
unaccounted-for cost of the fictional investigation.

SVU repays this much counting because of its unusual position in
American media. Few scripted programs have run long enough to constitute
the stable, decades-long message system that cultivation theory was
built to study. Fewer still occupy a topic where public beliefs bear so
directly on institutional behavior: how sex-crime allegations should be
investigated, how much collateral suspicion is tolerable, and what the
police owe the people they wrongly pursue. A viewer who has watched even
a fraction of the catalogue encounters hundreds of worked examples of
those questions, each resolved the same way. This article documents what
the worked examples contain.

\subsection{2. The Procedural and Its
Casualties}\label{the-procedural-and-its-casualties}

\subsubsection{2.1 Cultivation, crime drama, and what viewers take
away}\label{cultivation-crime-drama-and-what-viewers-take-away}

Cultivation research holds that heavy television viewing pulls beliefs
about the world toward the world television depicts (Gerbner et al.,
2002). Crime programming is the classic site for the effect, because the
genre's distortions are large and consistent: violent crime is
overrepresented, clearance is near-certain, and investigation is fast.
The survey record is messier than the theory. Dowler (2003) found
crime-drama viewing weakly linked to fear of crime and, in his own
conclusion, unrelated to perceived police effectiveness. Kort-Butler and
Sittner Hartshorn (2011) found that crime dramas predicted support for
the death penalty through a path that did not run through fear, while
nonfiction crime programming eroded confidence in the justice system by
stoking fear. Surette (2015) frames the underlying mechanism as a
substitution: for most viewers, mediated experience of criminal justice
is nearly all the experience they have.

The CSI-effect literature concedes on all sides that procedural
television plausibly shapes expectations about how investigation and
evidence work (Cole \& Dioso-Villa, 2009; Podlas, 2006). If jurors can
absorb a forensic fantasy from CSI, viewers can absorb an
interrogation-room ethics from SVU, and content analysis is the tool for
reading what that ethics contains (Eschholz et al., 2004).

The content-analytic tradition in this area has usually worked from
samples: a season here, a composite week of programming there, coded by
hand because hand-coding is expensive. Eschholz et al.~(2004) built
their analysis of NYPD Blue and Law \& Order from single-season samples.
A sample is a compromise, though, and for questions about a show's
conventions the compromise bites, because conventions are what varies
least within a season and most across eras. The present study codes the
population: every episode the program has aired. Whatever else may be
said about the instrument that made this affordable, the census design
removes an entire class of inferential worry.

\subsubsection{2.2 Copaganda and the SVU
literature}\label{copaganda-and-the-svu-literature}

A second body of work reads the procedural as advocacy. Kappeler and
Potter (2018) catalogue the mythologies that crime media sustain.
Karakatsanis (2025) popularized the term copaganda for the phenomenon;
his target is news coverage, though the entertainment genre raises the
same questions. Rapping (2003) reads the Law \& Order franchise
specifically as a long argument for prosecution as the solution to
social disorder.

SVU and its genre have drawn their own scholarship, most of it centered
on victims. Cuklanz (2000) traces how prime-time rape narratives evolved
from the 1970s onward, with the detective's sensitivity becoming the
genre's proof of enlightenment; Projansky (2001) maps how pervasively
rape structures postfeminist film and television. On SVU itself, Britto
et al.~(2007) find the show's victim pool skews young, white, and female
relative to crime statistics, and Cuklanz and Moorti (2006) credit the
series with a prime-time feminism that takes sexual violence seriously
while noting the limits of its politics. Bernabo (2022) examines how
broadcast procedurals answered the post-2020 reckoning over policing,
finding varied narrative strategies rather than a single reform posture.
What the SVU literature has needed is a systematic account of the show's
other recurring character: the innocent suspect. Critics of the genre
have long observed that procedurals churn through wrong men on the way
to the right one. Nobody, to my knowledge, has counted them, graded
their injuries, and checked who apologized.

That last item matters more than it may appear, because acknowledgment
of error is central to institutional legitimacy (Tyler, 2006); a drama
that shows harm without acknowledgment is modeling an accountability
norm, hundreds of times.

One conceptual point the copaganda framing sometimes flattens:
pro-police messaging need not hide harm; it can display harm and detach
it from responsibility. A series that showed no wrongful arrests would
be implausible to its own audience. A series that shows wrongful arrests
constantly, attaches them to sympathetic protagonists, resolves them
within the hour, and moves on without consequence is doing something
subtler than concealment. Whether SVU fits that description is an
empirical question about its scripts.

\subsubsection{2.3 Wrongful accusation off
screen}\label{wrongful-accusation-off-screen}

Off screen, wrongful accusation is neither rare nor cheap, but it is
rarer than SVU implies and far more expensive. Garrett (2011) documents
how ordinary investigative practice produced the first 250 DNA
exonerations. Gross et al.~(2005) catalogue exonerations from 1989 to
2003, a count the National Registry of Exonerations (2026) continues.
Loeffler et al.~(2019) estimate a wrongful conviction rate near 6\% from
prisoner self-reports. The aftermath is long: lasting psychiatric injury
(Grounds, 2004), years spent rebuilding identity and community
(Westervelt \& Cook, 2012). And the research consensus places knowingly
false reports of sexual assault at 2\% to 10\% of reports (Lisak et al.,
2010), with the stigma of a sex-crime allegation attaching regardless of
outcome. Sex crimes are SVU's entire jurisdiction.

\subsection{3. Method}\label{method}

\subsubsection{3.1 Sample and the state of the
source}\label{sample-and-the-state-of-the-source}

The sample is a census: all 576 episodes of SVU, seasons 1 through 27,
with season 27 partial at three episodes because it was airing at
collection time. Full episode transcripts were collected from
subslikescript.com and assembled into a single spreadsheet, one row per
episode.

That spreadsheet turned out to be corrupted in ways the coding
inherited, and the discovery is worth telling in order. The pipeline's
working assumption was simple: row i holds episode i's transcript.
Nothing in the coding output contradicted it. The failure surfaced only
when I audited the coded content against ground truth, comparing every
transcript's full text against every other's. Nine episodes' cells
contain a duplicate of another episode's script: S01E13 holds S01E15's,
and S02E11, S02E18, S04E10, S05E04, S05E14, S08E11, S09E19, and S11E09
likewise hold other episodes' transcripts (the full mapping is in the
supplementary material). The instrument coded the nine faithfully; none
drew its needs-deep-review flag, because each was a perfectly coherent
transcript of the wrong episode. The nine are excluded from every
episode-level denominator in this article, which is why the primary base
is 567 rather than 576, and the 15 person-records they generated are
removed from the person-level set. The assumption failed silently; only
ground-truth verification caught it.

The same audit quantified a second defect: 145 of the 576 transcripts
(25.2\%) end at 32,767 characters to the character, the Excel cell cap,
their final scenes cut off. The truncation is lopsided, 138 of the 145
in seasons 1-18 (33.7\% of that era) against 2 and 5 in the later eras,
and it censors endings, where resolutions, apologies, and final-scene
harms live. For prevalence and harm the direction is conservative: a
lost ending can only hide qualifying persons, injuries, and deaths, so
those counts are floors, the death audit included. For apology the same
missingness cuts the other way. An apology delivered in a lost final
scene is an apology the instrument never saw, so the 92.1\% no-apology
rate may be partly inflated, with the extra uncertainty concentrated in
the early era where the truncation sits. A repair of the 145 truncated
transcripts and a re-code of those episodes is planned and will settle
the question. What survives either way is the within-dataset gradient of
Section 4.7, apology tracking severity and ignoring conduct: the
gradient compares groups inside the same corpus, so endings lost without
regard to content thin every cell alike and leave the comparison
standing, a protection the marginal rate does not have.

Five further episodes had defective text of more ordinary kinds: four
(S14E01, S14E02, S15E01, S17E01) had empty transcripts, and one (S06E09)
contained non-English character encoding. In total, then, 14 of 576
source episodes were defective. Sections 3.4 and 6 describe how each
kind was handled.

\subsubsection{3.2 The coding instrument}\label{the-coding-instrument}

Every transcript was coded by a large language model,
claude-sonnet-4-5-20250929 (Anthropic), accessed through the Messages
Batch API. Each of the 576 requests contained the same fixed system
prompt, which is the complete codebook, and one episode transcript. The
model returned structured JSON: episode-level flags plus one record per
qualifying person. All 576 requests returned output without error, at a
cost between \$15 and \$17 in API fees and roughly 14 hours of waiting.

I treated the model as an instrument rather than an assistant. It never
saw my hypotheses and never participated in a conversation; it exercised
no discretion beyond what the codebook grants any coder: one transcript
in, one fixed rubric applied, one structured record out, 576 times. It
sits closer to an optical scanner than to a research collaborator. Like
any instrument it has error, and that error is described here but not
yet quantified: a stratified human validation study, prioritizing the
death rows in Section 4.8, is in progress. The design sits inside an
emerging methods literature on language models as annotators, which
insists on downstream validation (Gilardi et al., 2023; Törnberg, 2023;
Ziems et al., 2024), measured against classical content-analytic
reliability standards (Krippendorff, 2019; Neuendorf, 2017).

The batch design was chosen for discipline as much as cost. Because all
576 requests were submitted as a single job against a frozen prompt, no
coding decision could contaminate a later one. Batch independence rules
out sequential drift, the instrument learning my preferences or being
nudged mid-run; it does not rule out run-to-run variance, which is one
of the validation study's jobs. The output pipeline was mundane, and I
record it for replication. The model returned each episode's coding as
JSON; despite instructions it wrapped most responses in markdown fences,
stripped in post-processing. Records were then aggregated into two flat
files, one row per episode and one row per qualifying person, and those
two files are the sole source of every statistic in this article. The
full codebook accompanies the submission as supplementary material, and
the batch input can be regenerated from the transcript file, so the
measurement is repeatable by anyone for the price of a modest dinner.

Affordability is the methodological point: hand-coding 576 hour-long
episodes on a twenty-field codebook is a multi-year project or a funded
team, and this coding cost less than \$20. The trade is that
reliability, which a team of human coders establishes as it works, must
here be established afterward, by audit. Section 3.5 describes the
audits already performed; the Limitations section describes the one
still owed.

\subsubsection{3.3 Codebook and inclusion
criteria}\label{codebook-and-inclusion-criteria}

The codebook's inclusion rule, which I refer to as the North Star
criteria, admits a person-record only if all three conditions hold: (1)
the person is proven or strongly implied innocent of the accused
offense; (2) the person experiences public exposure (workplace, family,
school, media, church, or online) or a physical or violent consequence;
and (3) the harm follows from the accusation itself rather than from
unrelated events. Being accused in front of colleagues with no further
sanction qualifies, at the low end of the severity scale. Private
questioning with no spread and no material harm does not qualify at all.

Each record carries controlled-vocabulary fields: innocence status
(proven innocent, strongly implied innocent, partially involved,
actually guilty); role in plot (initial suspect, red herring, family
member, colleague, community member, other); the offense accused of;
accusation origin, defined as who first pointed the finger (victim
identification, witness misidentification, squad inference, coerced
interview, technical or database error, fabrication, unknown); exposure
channel and the party who disclosed; consequence category and free-text
detail; and a four-point severity scale: 1, private or low-level harm;
2, public exposure without formal sanction; 3, material sanction or
serious fallout; 4, life-altering outcomes or death. Two further fields
code police behavior: conduct toward the person (none, verbal threat,
coercive tactic, insult or degradation, multiple) and apology (none,
partial, formal, unknown), where a partial apology is softened language
without responsibility and a formal apology is a clear on-screen
acknowledgment of the wrongful accusation. Records also store verbatim
supporting quotes, free-form tags, and notes.

Two design choices are deliberate rather than incidental. The North Star
criteria are conservative, built to exclude rather than include: a
person privately questioned and released, with no spread and no material
consequence, generates no record, however unpleasant the scene. This
buys precision at the price of recall, and it means the counts reported
here are floors. And the retention of partially involved persons follows
from the study's definition of its object. A man who lied about an
affair and was then wrongly accused of the rape is innocent of the
accusation; the harm the accusation causes him is the phenomenon under
measurement, whatever else he did. Excluding such cases would quietly
adopt the detectives' own logic, that prior bad acts make a person a
legitimate target, and that logic is part of what the study exists to
examine. The codebook carries one documented defect here: its rules line
says exclude the guilty while its vocabulary includes
partially\_involved, a conflict at the margin that Section 3.4 records
resolving.

\subsubsection{3.4 Analysis set, exclusions, and
normalization}\label{analysis-set-exclusions-and-normalization}

The raw coding produced 541 person-records. Five are coded actually
guilty: Thomas ``Bird'' Gordon and Judge Walter Thornburg (both S03E19),
Gabriel Duvall in his initial questioning (S06E20), Bill Harris
initially (S12E03), and Orville Underwood (S12E12). I excluded the five:
the study counts harm to the innocent, and the codebook itself instructs
their exclusion. The transcript audit then removed 15 more: the phantom
records generated by the nine wrong-transcript episodes, among them
Stephanie Mulroney's S01E13 row and Larry Tauber's S05E04 row, both
codings of another episode's script. The analysis set is therefore n =
521.

The unit is the person-episode, and I use that term deliberately. Eight
records beyond the phantoms belong to recurring persons accused in
separate episodes: Nick Amaro and Marsden are coded three times each,
and Brian Cassidy, Elliot Stabler, Ed Tucker, and Lomatin twice each.
The 521 person-episode records therefore describe roughly 513 unique
persons. Records coded partially involved (n = 43) are retained: by the
controlled-vocabulary definition, these people are guilty of something
else and innocent of the accused offense, which is the situation the
study measures.

Two fields required normalization, disclosed here in full. In accusation
origin, the coder produced four out-of-vocabulary records labeled
victim\_misID (S12E15, S16E13, S20E09, S21E05); each is a victim
personally identifying the wrong person, the documented definition of
victim\_ID, so the four were folded in (90 + 4 = 94). One record labeled
third\_party (S17E17) matches the documented definition of
witness\_misID and was folded in (38 + 1 = 39). In the accused-of field,
three spelling variants of kidnapping were unified, and five
undocumented values (murder, kidnapping, arson, terrorism, child
abandonment) were retained as coded. The apology field contains two
records coded unknown (both S20E21); these are treated as non-apologies
in every rate.

The four transcriptless episodes returned has\_false\_suspect = Maybe
and zero person-records. Because 15 episodes in total carry the Maybe
flag, every episode-level rate in this article is reported under both
treatments, Y-only and Y-plus-Maybe; the overall false-suspect rate
moves by 2.7 percentage points between them, and no conclusion depends
on the choice. The single encoding-damaged episode (S06E09) still
yielded one person-record, which I treat as lower confidence. Sixteen
episodes (2.8\%) carry the coder's needs-deep-review flag. Those 16 have
not yet been individually re-reviewed; four are the transcriptless
episodes, and the rest are queued for the validation pass.

\subsubsection{3.5 Verification, the death audit, and
quotes}\label{verification-the-death-audit-and-quotes}

The free-text consequence field is a keyword trap, and I fell into it
before I climbed out. An early internal draft of this analysis counted
deaths by keyword matching and reported 31 murders, 21 suicides, and 25
assaults. Row-level verification against the consequence detail, notes,
and quoted scenes showed all three figures were wrong: the keyword count
had absorbed strings such as arrested\_for\_wrong\_murder,
suicide\_watch, suicide\_attempt, assisted\_suicide, family\_murdered,
and suspected\_then\_cleared\_by\_death, none of which is a death of the
accused person. Every death and injury figure in Section 4.8 was
therefore adjudicated row by row: each qualifying record inspected, each
borderline ruled on by name, each keyword trap excluded. Three injuries
that appear only in the notes field, never in the structured consequence
detail, are reported separately and never merged into headline counts. I
flagged this correction in the first draft because it taught a simple
rule: structured fields are for counting, free text is for reading.

The transcript audit extended that rule one level down. Structured
fields are only as trustworthy as the source cell they were computed
from, and the audit found deaths hiding in both directions: a murder
recorded only in the free-form tags (Dr.~Archibald Newlands, S05E05),
which the first death pass missed, and a suicide duplicated across two
episodes by a copied transcript (Larry Tauber), which the first pass
counted twice. The rules were tightened accordingly, and Section 4.8
reports the adjudicated final numbers with every borderline ruling
named.

All quoted dialogue in this article was verified verbatim against the
source transcripts. The transcripts carry no speaker labels, so where a
quote is attributed to a named character, the attribution follows scene
context, and the same fragility runs through the instrument itself.
Three coded fields depend directly on knowing who is speaking: police
conduct (was the threat delivered by a detective), the disclosing party
in exposure (who told the employer or convened the cameras), and police
apology. Every value in those fields rests on the model inferring
speakers from context, and an inference that systematically mistook,
say, an ADA's threats for a detective's would move the rates while
leaving no other trace. The most exposed number in this article is the
83.9\% squad share of disclosure, which aggregates hundreds of such
attributions. The validation study therefore includes speaker
attribution as a named task: human raters will re-attribute the
supporting quotes for a stratified sample of conduct, disclosure, and
apology records against video or labeled scripts.

\subsubsection{3.6 Analytic approach}\label{analytic-approach}

Because the dataset is a census of the show rather than a sample, the
results are population descriptions of the text and carry no sampling
error; I report proportions and cross-tabulations without significance
tests. The companion longitudinal study tests era contrasts, under an
explicitly stated view of each season as one draw from a stochastic
writers'-room process (Orden, in preparation); for the census
description here, the proportions are the facts. Severity means treat
the four-point scale as interval for summary only, and medians are
reported alongside. All figures were computed directly from the two flat
files with a Python standard-library script, and I validated every
number against them.

\subsection{4. Results}\label{results}

\subsubsection{4.1 Prevalence}\label{prevalence}

Table 1 reports the episode-level rates on the corrected base. Of 567
episodes, 337 (59.4\%) contain at least one person accused and later
shown or strongly implied to be innocent; adding the 15 ambiguous
episodes raises the figure to 352 (62.1\%). Public exposure of an
innocent person occurs in 295 episodes (52.0\%; 306 or 54.0\% with
ambiguous cases). More than half the time, an hour of SVU includes the
public shaming of somebody who did not do it.

\textbf{Table 1. Episode-level prevalence (N = 567 episodes)}

{\def\LTcaptype{none} % do not increment counter
\begin{longtable}[]{@{}
  >{\raggedright\arraybackslash}p{(\linewidth - 8\tabcolsep) * \real{0.2000}}
  >{\raggedright\arraybackslash}p{(\linewidth - 8\tabcolsep) * \real{0.2000}}
  >{\raggedright\arraybackslash}p{(\linewidth - 8\tabcolsep) * \real{0.2000}}
  >{\raggedright\arraybackslash}p{(\linewidth - 8\tabcolsep) * \real{0.2000}}
  >{\raggedright\arraybackslash}p{(\linewidth - 8\tabcolsep) * \real{0.2000}}@{}}
\toprule\noalign{}
\begin{minipage}[b]{\linewidth}\raggedright
Flag
\end{minipage} & \begin{minipage}[b]{\linewidth}\raggedright
Y
\end{minipage} & \begin{minipage}[b]{\linewidth}\raggedright
Maybe
\end{minipage} & \begin{minipage}[b]{\linewidth}\raggedright
Y-only rate
\end{minipage} & \begin{minipage}[b]{\linewidth}\raggedright
Y+Maybe rate
\end{minipage} \\
\midrule\noalign{}
\endhead
\bottomrule\noalign{}
\endlastfoot
Contains a false suspect & 337 & 15 & 337/567 = 59.4\% & 352/567 =
62.1\% \\
Contains public exposure of an innocent person & 295 & 11 & 295/567 =
52.0\% & 306/567 = 54.0\% \\
Flagged for deep review & 16 & n/a & 16/567 = 2.8\% & n/a \\
\end{longtable}
}

At the person level, the 521 qualifying records work out to 0.92 per
episode: nearly one innocent life disrupted per hour of television,
sustained for 27 years. Table 2 gives the per-season counts. The early
and middle seasons run hottest, with season 9 the densest at 1.89
persons per episode, and the rate falls in the 2020s, bottoming at 6
persons in season 21 (0.30 per episode), though it never reaches zero.

\textbf{Table 2. Persons harmed per season (episode counts exclude
corrupted-source episodes)}

{\def\LTcaptype{none} % do not increment counter
\begin{longtable}[]{@{}llll@{}}
\toprule\noalign{}
Season & Episodes & Persons harmed & Persons/episode \\
\midrule\noalign{}
\endhead
\bottomrule\noalign{}
\endlastfoot
1 & 21 & 17 & 0.81 \\
2 & 19 & 29 & 1.53 \\
3 & 23 & 25 & 1.09 \\
4 & 24 & 32 & 1.33 \\
5 & 23 & 34 & 1.48 \\
6 & 23 & 27 & 1.17 \\
7 & 22 & 18 & 0.82 \\
8 & 21 & 22 & 1.05 \\
9 & 18 & 34 & 1.89 \\
10 & 22 & 27 & 1.23 \\
11 & 23 & 29 & 1.26 \\
12 & 24 & 22 & 0.92 \\
13 & 23 & 29 & 1.26 \\
14 & 24 & 20 & 0.83 \\
15 & 24 & 15 & 0.62 \\
16 & 23 & 14 & 0.61 \\
17 & 23 & 16 & 0.70 \\
18 & 21 & 17 & 0.81 \\
19 & 24 & 15 & 0.62 \\
20 & 24 & 18 & 0.75 \\
21 & 20 & 6 & 0.30 \\
22 & 16 & 8 & 0.50 \\
23 & 22 & 9 & 0.41 \\
24 & 22 & 13 & 0.59 \\
25 & 13 & 8 & 0.62 \\
26 & 22 & 14 & 0.64 \\
27 & 3 & 3 & 1.00 \\
\textbf{Total} & \textbf{567} & \textbf{521} & \textbf{0.92} \\
\end{longtable}
}

\subsubsection{4.2 Who gets accused, and of
what}\label{who-gets-accused-and-of-what}

Innocence is rarely left ambiguous. Of the 521 persons, 429 (82.3\%) are
proven innocent on screen through DNA, a confession by the real
perpetrator, or an explicit exoneration; 49 (9.4\%) are strongly implied
innocent; 43 (8.3\%) are partially involved, guilty of some lesser wrong
but innocent of the accusation. The show wants the audience to know
these people did not do it.

The accusation itself is usually rape (316 persons, 60.7\%), followed by
assault (65, 12.5\%), child sexual abuse (45, 8.6\%), and vaguely
specified sex crimes (39, 7.5\%), consistent with the unit's
jurisdiction. Structurally, the accused fall into two large groups:
initial suspects, the first person the squad pursues (205, 39.3\%), and
red herrings written into the script to mislead the viewer (172,
33.0\%). The two groups are treated differently by the narrative.
Initial suspects carry a mean severity of 2.98 and reach severity 4 in
28.8\% of cases; red herrings average 2.44 and reach severity 4 in
13.4\%. Family members of victims or suspects (78, 15.0\%) sit in
between at a mean of 2.68, with colleagues (24), community members (19),
and others (23) making up the remainder.

The free-form tags sketch who these people are, with a caution: the tags
are the coder's own free vocabulary, unaudited, so the counts are
indicative only. On that footing, 51 persons are tagged as teenagers, 40
are entangled in custody disputes, 28 are police officers, 22 are
teachers, 13 are doctors, and 13 are immigrants: statuses in which
accusation alone ends careers or triggers deportation machinery.

\subsubsection{4.3 Severity}\label{severity}

Table 3 shows the distribution across the four-point scale, and the
finding is its center of gravity: 325 of 521 persons (62.4\%) land at
severity 3 or 4, with a mean of 2.72 and a median of 3. Only 60 persons
(11.5\%) escape with private, low-level harm.

\textbf{Table 3. Consequence severity (n = 521 persons)}

{\def\LTcaptype{none} % do not increment counter
\begin{longtable}[]{@{}llll@{}}
\toprule\noalign{}
Severity & Definition & n & \% \\
\midrule\noalign{}
\endhead
\bottomrule\noalign{}
\endlastfoot
1 & Private or low-level harm & 60 & 11.5\% \\
2 & Public exposure, no formal sanction & 136 & 26.1\% \\
3 & Material sanction or serious fallout & 215 & 41.3\% \\
4 & Life-altering outcome or death & 110 & 21.1\% \\
\end{longtable}
}

Most persons are hurt in more than one domain at once: 204 (39.2\%) are
coded into multiple consequence categories, with legal consequences
alone accounting for a further 152 (29.2\%), and social, work, physical,
and family harm dividing the rest.

\subsubsection{4.4 Where accusations come
from}\label{where-accusations-come-from}

Table 4 reports the accusation origin, the controlled field recording
who first pointed the finger, and its lead finding: the single largest
source of false accusation, at 284 of 521 persons (54.5\%), is squad
inference, the detectives' own theory, built from pattern-matching,
proximity, or a prior record. Some of that share is definitional. A
detective procedural runs on detectives generating suspects, so the
squad leading its own origin table is partly the genre describing
itself, and I read the field as context, locating where error enters the
story, rather than as an indictment on its own; the indictment sits in
the fields that measure what happens next, in severity, disclosure, and
apology. What keeps the origin finding above plot mechanics is the
inclusion rule itself: each of these 284 inferences produced public
exposure or violence for an innocent person.

\textbf{Table 4. Accusation origin (n = 521 persons)}

{\def\LTcaptype{none} % do not increment counter
\begin{longtable}[]{@{}lllll@{}}
\toprule\noalign{}
Origin & n & \% & Mean severity & \% severity 4 \\
\midrule\noalign{}
\endhead
\bottomrule\noalign{}
\endlastfoot
Squad inference & 284 & 54.5\% & 2.50 & 14.1\% \\
Victim identification & 94 & 18.0\% & 2.99 & 25.5\% \\
Fabrication & 66 & 12.7\% & 3.24 & 42.4\% \\
Witness misidentification & 39 & 7.5\% & 2.44 & 7.7\% \\
Technical/database error & 25 & 4.8\% & 2.92 & 32.0\% \\
Coerced interview & 12 & 2.3\% & 3.33 & 50.0\% \\
Unknown & 1 & 0.2\% & 4.00 & 100.0\% \\
\end{longtable}
}

One reconciliation is needed, because the free-form tags look like a
different story. The tag wrong\_ID appears on 309 of the innocent
records and could be misread as evidence that misidentification is the
leading cause. It is a different measurement: the tag marks an episode
outcome, ``the squad had the wrong person,'' and among the 309 tagged
records the controlled origin field is squad inference in 185 (59.9\%),
victim identification in 67, and witness misidentification in only 29.
The controlled field records who first pointed the finger, and on that
field the squad's own reasoning dominates.

The severity columns of Table 4 carry a second finding. Fabrication and
coerced interviews, the origins involving deliberate wrongdoing, produce
the most catastrophic outcomes (42.4\% and 50.0\% severity 4, the latter
on a small n of 12), while squad inference, the most common origin,
produces the gentlest (14.1\% severity 4, mean 2.50). I return to this
asymmetry in Section 5.4.

\subsubsection{4.5 Exposure and the media
multiplier}\label{exposure-and-the-media-multiplier}

For 377 of 521 persons (72.4\%), the accusation spreads beyond law
enforcement. The two largest channels, workplace (143 persons, 27.4\%)
and police-only containment (142, 27.3\%), are a single record apart, a
difference on which nothing should be built; family (90, 17.3\%) and
media (86, 16.5\%) follow, with the rest divided among schools, courts,
churches, online, and multiple channels.

Who spreads the word is less ambiguous. The squad is coded as the
disclosing party for 454 of the 521 persons (87.1\%), but that raw
figure includes 136 cases that never left the precinct. The fairer
version is conditional: among the 379 persons whose accusation traveled
beyond the police, detectives are the ones who told the employer,
confronted the person in public, or staged the arrest for an audience in
318 cases (83.9\%), which is 61.0\% of all persons in the dataset, a
figure that leans on inferred speaker attribution in unlabeled
transcripts more heavily than any other in this article (Section 3.5).
Victims disclose in 4.8\% of cases, and the media originate the exposure
in only 3.3\%.

Where the exposure lands tracks how bad things get. Table 5 crosses
channel with severity: cases contained to the police average 2.04, with
8.5\% reaching severity 4, while media exposure averages 3.31, with
45.3\% reaching severity 4, a 5.4-fold difference in catastrophic
outcomes. The small legal-exposure category runs even higher (57.1\%
severity 4 on n = 14). A caution before reading this causally: channel
and severity are co-authored by the same scripts, and an episode built
as a catastrophe is also one likely to feature cameras, so narrative
prominence confounds any claim that publicity produces the damage. What
the table establishes is dramaturgical: the show pairs media exposure
with catastrophe, consistently, across 27 seasons.

\textbf{Table 5. Severity by exposure channel (n = 521 persons)}

{\def\LTcaptype{none} % do not increment counter
\begin{longtable}[]{@{}lllll@{}}
\toprule\noalign{}
Channel & n & \% of persons & Mean severity & \% severity 4 \\
\midrule\noalign{}
\endhead
\bottomrule\noalign{}
\endlastfoot
Workplace & 143 & 27.4\% & 2.78 & 15.4\% \\
Police only & 142 & 27.3\% & 2.04 & 8.5\% \\
Family & 90 & 17.3\% & 2.84 & 20.0\% \\
Media & 86 & 16.5\% & 3.31 & 45.3\% \\
School & 20 & 3.8\% & 3.00 & 20.0\% \\
Multiple & 18 & 3.5\% & 3.33 & 33.3\% \\
Legal & 14 & 2.7\% & 3.57 & 57.1\% \\
Church & 4 & 0.8\% & 2.75 & 25.0\% \\
Online & 2 & 0.4\% & 3.00 & 0.0\% \\
Unknown & 2 & 0.4\% & 1.00 & 0.0\% \\
\end{longtable}
}

Read together, the two findings are damning on the show's own terms:
publicity accompanies the worst outcomes, and in 83.9\% of the cases
where word spreads, the squad is the one spreading it.

\subsubsection{4.6 How detectives treat the
innocent}\label{how-detectives-treat-the-innocent}

Just over half of the innocent persons in the dataset, 276 of 521
(53.0\%), are threatened, coerced, or degraded by detectives at some
point in their episode. Table 6 gives the breakdown; verbal threats
lead, with coercive tactics such as fabricated evidence claims and false
promises close behind.

\textbf{Table 6. Police conduct toward the innocent person (n = 521)}

{\def\LTcaptype{none} % do not increment counter
\begin{longtable}[]{@{}lll@{}}
\toprule\noalign{}
Conduct & n & \% \\
\midrule\noalign{}
\endhead
\bottomrule\noalign{}
\endlastfoot
None documented & 245 & 47.0\% \\
Verbal threat & 108 & 20.7\% \\
Coercive tactic & 68 & 13.1\% \\
Multiple types & 55 & 10.6\% \\
Insult or degradation & 45 & 8.6\% \\
\textbf{Any threat, coercion, or degradation} & \textbf{276} &
\textbf{53.0\%} \\
\end{longtable}
}

Two scenes, a quarter century apart, put flesh on those categories. In
S01E12, detectives threaten a sex worker named Katya Ivanova with
deportation: ``we know you're illegal\ldots{} if you're not at our
office by 9:00 a.m., Immigration will be at your house by 9:05. That I
can personally guarantee you.'' Twenty seasons later, in S21E06, Carlos
Hernandez describes the interrogation that took 16 years of his life:
``Because Monte and the DA told me if I didn't confess, that I'd get the
sayonara syringe. Man, I was 18. I didn't have no mommy no more to help
me.''

The dataset does not distinguish beloved series regulars from
one-episode bad cops; what it records is that the person on the
receiving end was innocent, and that this happens in more than half of
all qualifying cases.

\subsubsection{4.7 The apology gradient}\label{the-apology-gradient}

Acknowledgment is where the numbers get quiet. Across 521 innocent
persons, detectives offer no apology at all in 480 cases (92.1\%).
Partial apologies, softened language without responsibility, occur 31
times (6.0\%). Formal on-screen acknowledgment of the wrongful
accusation occurs 8 times: 1.5\% of cases, spread over 27 seasons (Table
7). The coding originally showed nine formal-apology rows, with
Stephanie Mulroney apparently apologized to in two consecutive episodes.
The audit dissolved the oddity: the S01E13 cell contains S01E15's
script, so the one apology in ``Entitled'' had been coded twice.

\textbf{Table 7. Police apology (n = 521 persons)}

{\def\LTcaptype{none} % do not increment counter
\begin{longtable}[]{@{}lll@{}}
\toprule\noalign{}
Apology & n & \% \\
\midrule\noalign{}
\endhead
\bottomrule\noalign{}
\endlastfoot
None & 480 & 92.1\% \\
Partial & 31 & 6.0\% \\
Formal & 8 & 1.5\% \\
Unknown & 2 & 0.4\% \\
\textbf{Any apology (partial + formal)} & \textbf{39} &
\textbf{7.5\%} \\
\end{longtable}
}

Who gets the 39 apologies is more revealing than how few there are.
Table 8 shows the severity gradient: none of the 60 severity-1 persons
receives even a partial apology; severity-2 and severity-3 persons
receive one 6.6\% and 5.6\% of the time; severity-4 persons, the
destroyed and the dead, receive one 16.4\% of the time, and five of the
eight formal apologies belong to them. Misconduct, by contrast, moves
the rate by nothing measurable: persons who were threatened or coerced
are apologized to at 8.0\%, persons who were treated correctly at 6.9\%.
Apology in this fiction tracks catastrophe, never conduct. Detectives
say sorry when someone is ruined or dead, at odds of about one in six;
they do not say sorry because of anything they themselves did.

\textbf{Table 8. Apology by severity and by police conduct (n = 521
persons)}

{\def\LTcaptype{none} % do not increment counter
\begin{longtable}[]{@{}lllll@{}}
\toprule\noalign{}
Group & n & Any apology & Rate & Formal \\
\midrule\noalign{}
\endhead
\bottomrule\noalign{}
\endlastfoot
Severity 1 & 60 & 0 & 0.0\% & 0 \\
Severity 2 & 136 & 9 & 6.6\% & 2 \\
Severity 3 & 215 & 12 & 5.6\% & 1 \\
Severity 4 & 110 & 18 & 16.4\% & 5 \\
Threatened or coerced & 276 & 22 & 8.0\% & n/a \\
Not threatened & 245 & 17 & 6.9\% & n/a \\
\end{longtable}
}

The eight formal apologies are the entire on-screen accountability
record of a 567-episode corpus, so they can be listed. Captain Cragen
tells Stephanie Mulroney (S01E15), ``I was wrong, Ms.~Mulroney,'' and,
over her flat ``All right,'' asks her to ``please accept my apologies to
your family.'' A court apologizes to Evan Ramsey (S03E01), who is shot
dead days later by his own daughter. A detective quietly corrects the
record for Varla, a homeless woman with a phantom pregnancy accused of
abandoning a baby (S07E14). Victor Tate (S11E01) hears Stabler say, ``I
made a mistake, Victor. I wanted to tell you myself''; Tate's answer is
quoted in Section 5.2. Frank Sullivan (S11E21) receives his from an ADA
after he slits his wrists in an interrogation room. Detectives reveal
that Vicki Harris's arrest for prostitution (S13E10) was staged to get
her out of danger. Omar Pena's exoneration (S13E17) carries two
acknowledgments, one private and one official: Benson admits to a
colleague, ``I broke him,'' and then, ``I forced a false confession,''
and it is the judge, vacating the conviction, who tells Pena in open
court that ``this court sincerely apologizes.'' Stephen Lomatin
(S17E09), an auxiliary police officer accused of rape, receives the
eighth. That is the complete list.

\subsubsection{4.8 The dead}\label{the-dead}

Table 9 summarizes the adjudicated death and physical-harm audit
described in Section 3.5. Every count below was verified record by
record, and every borderline case was ruled on by name.

\textbf{Table 9. Adjudicated deaths and physical harm among the 521
innocent persons}

{\def\LTcaptype{none} % do not increment counter
\begin{longtable}[]{@{}
  >{\raggedright\arraybackslash}p{(\linewidth - 2\tabcolsep) * \real{0.5000}}
  >{\raggedright\arraybackslash}p{(\linewidth - 2\tabcolsep) * \real{0.5000}}@{}}
\toprule\noalign{}
\begin{minipage}[b]{\linewidth}\raggedright
Category
\end{minipage} & \begin{minipage}[b]{\linewidth}\raggedright
Final count
\end{minipage} \\
\midrule\noalign{}
\endhead
\bottomrule\noalign{}
\endlastfoot
Murdered or killed as a consequence of the accusation & 21 (2 in
custody; 1 killed by police) \\
Completed suicides & 9 \\
Other accusation-linked deaths (overdose; AIDS after wrongful
imprisonment) & 2 \\
\textbf{Total dead} & \textbf{32, across 31 distinct episodes} \\
Apologies (partial or formal) among the dead & 5 of 32 \\
Suicide attempts, survived & 8, all in structured fields \\
Non-fatal assaults, vigilante attacks, and shootings & 28 in structured
fields, plus 3 verified in notes only \\
\end{longtable}
}

Twenty-one innocent persons are murdered or killed as a consequence of
the accusation. Two die in custody: Russell Ramsay, beaten, raped, and
killed at Rikers (S02E08), and Manny Montero, murdered at Rikers after
agreeing to testify (S15E13). One is killed by the police themselves:
Terrence Reynolds (S17E05). The rest die outside, mostly at the hands of
people who believed the accusation. Eric Byers (S10E02), a 17-year-old
stepbrother of the victim, is among them. So is Jerome Jones (S16E13),
seized on the street trying to return a woman's purse; she cries
``Please! He attacked me,'' and moments later, ``Then arrest him! He
raped me!'', and after security footage clears him, the victim's father
kidnaps, tortures, and murders him anyway. One murder entered the count
only at adjudication: Dr.~Archibald Newlands (S05E05), whose death sits
in the free-form tags and never in the structured consequence field,
which is why the first pass missed him; the rules now sweep the tags
too.

Nine more die by suicide. Mariel Plummer (S05E18), a child-services
worker, is arrested, fired, stripped of her pension, and pilloried in
the press before killing herself. Edwin Adelson (S12E02), framed with
planted files, perp-walked, and outed in the media, jumps in front of a
train after arraignment. Eddie Cappilla (S04E21) kills himself in the
same episode in which his brother Joe is murdered, making that hour the
only one in the series with two accusation-linked deaths. One suicide
moved houses during the audit: Larry Tauber's death had been coded to
S05E04, but that cell holds S06E14's script; the suicide is S06E14
content, where Tauber's genuine record already sat, and it is counted
once. Two deaths fall outside both categories. Leon Tate (S03E04) dies
of a drug overdose while wrongly suspected of serial murders, and his
mother blames the police. Ricardo Torres (S21E06), coerced into false
testimony at 15, spends 15 years in prison, contracts HIV there, and
dies of AIDS shortly after release.

In total, 32 people innocent of the accusation against them die, across
31 distinct episodes: 5.4\% of the 576 aired hours, 5.5\% of the
corrected base. Among the 32 dead, police apologize, formally or
partially, in 5 cases.

Below the fatal line, 8 persons attempt suicide and survive, all
recorded in structured fields: among them Charlotte Truex (S08E14), a
minor who slits her wrist immediately after a coerced confession taken
without counsel, and Kevin Wright (S12E21), whose attempt sits in the
structured tags, where an earlier pass had misfiled it as notes-only.
Five named records were considered and excluded (Robert Logan S05E03,
Colin Bennett S18E15, Aidan Connor S05E15, Javier Vega S05E21, Ashlee
Walker S10E19): in each, the record supports ideation, threat, or
ambiguity rather than an attempt. Twenty-eight persons suffer non-fatal
violence recorded in structured fields, with three more verified only in
notes. Police gunfire accounts for five: Willie Angel (S04E13) is shot,
paralyzed, framed with a drop gun, and wrongfully imprisoned for 25
years; Yusef Barre (S15E11) is shot and paralyzed by a detective who
believed he was armed; Michael Baxter (S05E11) and Larry Luft (S11E13)
survive their wounds; and Xavier Gonzalez (S22E07) is fired upon and
missed. The remainder are vigilante beatings, family violence, and
in-custody assaults: Coach Alexei Belyakov (S14E13), stabbed with an ice
pick over an abuse accusation that proves false; Eric Plummer (S03E02),
wrongfully convicted and then beaten, raped, and tortured across seven
years in prison; Sean Roberts (S18E02), whose 16 years of wrongful
imprisonment included repeated rape, adjudicated in as in-custody
violence; and Luna Prasada (S21E15), a framed executive jailed without
bail amid a media pile-on and sexually assaulted by fellow inmates. One
of the three notes-only cases, Marc Rajic (S13E09), survives an attempt
to burn him alive with gasoline while the structured field records only
a vigilante threat.

The audit's exclusions define the counts as much as its inclusions, so I
list the traps by name. Travis Hall (S01E04) is tagged
suspected\_then\_cleared\_by\_death, but he overdosed a week before the
crime; the death precedes the accusation. Officer Peter Ridley (S01E04)
is arrested\_for\_wrong\_murder, a phrase containing the word murder and
no dead accused person. Carlos Hernandez (S21E06) carries a tag
recording that his mother and sister were murdered while he sat wrongly
imprisoned; he survived, so he counts toward no death figure, however
much he belongs in any moral accounting. Doug Loveless (S12E19) and Jean
Dussault (S01E08) carry death-keyword tags referring to other people's
deaths and a pre-accusation overdose. Dr.~Matt Spevak (S05E23) was
murdered by his twin and framed posthumously: death before accusation,
again. Mr.~Lee (S18E17), a restaurant owner, receives death threats and
an armed intrusion but no on-screen injury. And Guarana (S15E17) is
excluded from the audit entirely as a dataset inclusion error: she was
never accused, and her row is queued for removal in the validation pass.
Every one of these strings defeats a keyword counter, which is how an
earlier draft came to report 31 murders and 21 suicides. The adjudicated
figures are 21 and 9.

\subsection{5. Discussion}\label{discussion}

\subsubsection{5.1 The wrong man as
furniture}\label{the-wrong-man-as-furniture}

The prevalence numbers settle one question: the falsely accused person
is a load-bearing element of the SVU formula, present in three of every
five episodes and averaging nearly one qualifying victim per hour across
27 years. Red herrings alone account for 172 persons, characters the
writers created specifically to be wrongly suspected, whose narrative
function is to absorb suspicion for two acts and release it in the
third.

What the counting adds to that familiar observation is the injury
report. These disposable characters are damaged severely: 62.4\% at
severity 3 or 4, arrest and firing and divorce as standard equipment,
death for 32 of them. The convention, named plainly, is that serious
harm to the innocent is a normal, expected, and narratively unremarkable
byproduct of investigation. The show does show the harm; that is what
makes the dataset possible. What it declines to show, 92.1\% of the
time, is anyone answering for it.

Table 2 does record a real decline in volume, from 1.53 persons per
episode in season 2 down to 0.30 in season 21, and the companion study
takes that trajectory across the show's cultural eras in detail (Orden,
in preparation). For present purposes two facts suffice: the convention
never stops, even the quietest season injures someone, and the harm to
those who still qualify has grown no milder, with a suggestive spike
toward worse during the \#MeToo-era seasons (see the companion).

The severity gap between initial suspects (mean 2.98) and red herrings
(2.44) has its own dramaturgical logic. Red herrings exist for the
audience, so the script releases them cheaply at the reveal; initial
suspects exist for the detectives, and the story extracts full value:
the arrest, the arraignment, the family fracturing on the courthouse
steps. On the numbers, the suspects the squad believes in hardest are
the ones the fiction hurts worst.

\subsubsection{5.2 An economy of apology}\label{an-economy-of-apology}

The apology findings are, to me, the most legible single result in the
study, and they were the first thing in the data that genuinely
surprised me. I expected apologies to be rare. I did not expect them to
be indifferent to conduct. A detective who threatens an innocent
teenager and a detective who behaves impeccably apologize at nearly
identical rates, 8.0\% against 6.9\%. The only variable that moves
apology is the size of the wreckage, and even catastrophe only lifts the
rate to 16.4\%. Five of the eight formal apologies in the whole series
attach to severity-4 cases, and most of the dead still receive none:
police apologize in only 5 of the 32 fatal cases.

Procedural-justice research treats acknowledgment of error as a core
input to institutional legitimacy (Tyler, 2006). SVU models the opposite
equilibrium, several hundred times: error without acknowledgment,
followed by the next case. The lesson on offer is that when police harm
the innocent, apology is a discretionary grace triggered by ruin, never
an obligation triggered by conduct. The show's most sympathetic
characters administer this equilibrium, and that is what gives it force;
the norm is modeled by heroes.

The one scene that breaks the pattern proves it. When Stabler apologizes
to Victor Tate in S11E01, the script gives Tate the reply almost no one
else in the dataset gets to deliver: ``And then my life just stopped!
Lost my woman, my friends! Hell, my mother barely knows me!'' The show
knows what an accusation costs and can write the bill when it chooses.
It chooses eight times in 27 seasons. The CSI-effect literature spent a
decade asking whether television teaches jurors to expect too much of
forensic evidence (Cole \& Dioso-Villa, 2009; Podlas, 2006); the
parallel question raised here is whether television teaches audiences to
expect nothing at all when the state gets it wrong.

\subsubsection{5.3 Publicity as
punishment}\label{publicity-as-punishment}

SVU has an implicit theory of stigma, and the severity-by-channel table
recovers it with numbers. Accusation contained to the precinct is
survivable: mean severity 2.04, catastrophic outcomes at 8.5\%.
Accusation that reaches a camera is close to a sentence: 45.3\% of
media-exposed persons reach severity 4, the 5.4-fold difference at the
center of this study. The show understands, as Goffman (1963) did, that
the injury of stigma is administered by an audience, and it stages the
principle again and again: the perp walk, the reporters on the
courthouse steps, the neighbor watching the arrest. Within the fiction
this is dramaturgy rather than demonstrated causation; the scripts that
convene the cameras are the scripts that plan the catastrophe.

The uncomfortable half of the finding is who convenes the audience. In
83.9\% of the cases where the accusation spreads at all, the disclosing
party is the squad itself; the press originates exposure in 3.3\%.
Within the fiction, the institution that transforms a survivable
suspicion into a public catastrophe is the police department the series
celebrates, and the transformation is depicted as procedure: telling the
employer, confronting the suspect at work, staging the arrest. A show
built on sympathy for victims has spent 27 seasons scripting its own
detectives' communication habits as the most reliable amplifier of harm
to the innocent.

The ordinary case tells the same story at lower volume. Accusations
spread beyond the police for 72.4\% of persons, and the workplace, one
of the two largest channels at 27.4\%, is where the show's middle-class
professionals are questioned in front of colleagues and escorted past
their desks. The codebook's severity floor captures the logic: public
accusation with no formal sanction at all is already severity 2, because
the audience is the sanction. Keeping suspicion inside the building
reads on screen as an act of unusual restraint; it occurs in just over a
quarter of cases.

\subsubsection{5.4 Fiction against the
record}\label{fiction-against-the-record}

Setting the show's world beside the documented one requires care about
units. The best available estimate of the wrongful conviction rate, from
prisoner self-reports, sits near 6\% (Loeffler et al., 2019), and
exoneration registries count the confirmed cases in the low thousands
over three and a half decades (Gross et al., 2005; National Registry of
Exonerations, 2026). SVU stages a false accusation in 59.4\% of
episodes. An episode is a case, roughly, and a false suspect is upstream
of a wrongful conviction, so the two rates measure different points in
the funnel; no adjustment I can construct, however, closes a gap of that
order. The show's fictional world is one in which wrongful suspicion is
the majority experience of contact with a sex-crimes unit.

The distortion runs in a second direction with the opposite sign. On
SVU, the error is discovered within the hour, the real perpetrator
confesses, and innocence is proven on screen for 82.3\% of the harmed.
In the documented record, discovery takes years or decades, arrives for
a fraction of the wrongly convicted, and leaves psychiatric injury that
outlasts the exoneration (Garrett, 2011; Grounds, 2004; Westervelt \&
Cook, 2012). The series therefore rehearses two miscalibrations at once:
wrongful accusation as far more common than reality, and far more
self-correcting. A viewer who took both lessons on board would arrive at
something like the show's own implicit position, that collateral damage
is ubiquitous, temporary, and nobody's fault.

A related caution: the 12.7\% fabrication share in Table 4 is a share of
harmed persons in a drama, and it cannot be read as a false-report rate;
the consensus for knowingly false sexual assault reports, 2\% to 10\%,
uses a different denominator entirely (Lisak et al., 2010). What the
origin table does support is a claim about dramaturgy. The show reserves
its worst outcomes for accusations rooted in civilian deceit (42.4\%
catastrophic) and its mildest for accusations rooted in detective
reasoning (14.1\%), even though detective reasoning generates the
majority of false accusations in the first place, a majority the genre's
structure partly guarantees. Institutional error is rendered frequent
but gentle; civilian malice is rendered rare but lethal. That allocation
protects the heroes twice over, and it is exactly the pattern the
copaganda literature would predict a police-centered drama to produce
(Kappeler \& Potter, 2018; Karakatsanis, 2025).

Whether viewers absorb any of this is a question this study cannot
answer and deliberately does not claim to. Content sets the ceiling on
cultivation; reception decides what passes through, and the audience
research cuts both ways: Hust et al.~(2015) found Law \& Order franchise
exposure associated with lower rape-myth acceptance and better consent
negotiation among college students. Reception can diverge from the
message system. What the findings here specify is the message system
itself, and they hand experimental researchers concrete stimuli: the
uncommented coerced confessions, the eight formal apologies, the perp
walks that precede suicides. Whether exposure to the convention shifts
beliefs about police accountability or juror expectations about
accusation and error is now a designable experiment.

\subsection{6. Limitations}\label{limitations}

Seven limitations bound these findings. First, and most important, the
coding instrument has no human reliability estimate yet. On the fields I
could check mechanically, the instrument violated instructions at a
measurable rate: it returned five actually-guilty person-records despite
a codebook prohibition, produced out-of-vocabulary origin values,
introduced undocumented accused-of values, and wrapped output in
markdown fences it was told to omit. Errors measured on checkable fields
imply nonzero error on uncheckable ones, severity judgments above all.
The stratified human validation study, prioritizing the 32 death records
and the severity-4 stratum, is therefore a requirement for confidence,
not a formality. The model rates its own confidence high on 506 of 521
records (97.1\%); that figure is uncalibrated and substitutes for
nothing.

Second, the source text was imperfect in ways that survive the audit.
The nine wrong-transcript episodes are excluded and their phantom
records removed, but 145 transcripts (25.2\%) are truncated at the Excel
cell limit, and the truncation concentrates in seasons 1-18 (33.7\% of
that era). For harm the resulting bias is conservative: a lost ending
can only hide qualifying persons and deaths, so those counts are floors.
For apology it is the opposite of conservative, because the missing
final scenes are exactly where apologies would appear; the 92.1\%
no-apology rate therefore carries extra uncertainty in the early era,
and the planned repair and re-code of the truncated episodes is the
cure. The severity-and-conduct gradient of Section 4.7 is the apology
finding least exposed here, since lost endings thin its cells together.
Four episodes had no transcript and contribute only ambiguity flags; one
had encoding damage and contributes a single lower-confidence record;
season 27 was three episodes old at collection.

Third, the coder read transcripts, never footage, so tone, framing, and
performance are invisible to it. A raised eyebrow can turn a threat into
a bluff both parties understand; no transcript records the eyebrow. A
single instrument also means any interpretive bias it carries is applied
uniformly across all episodes. Uniform bias leaves comparisons between
categories intact, the media multiplier and the apology gradient among
them; absolute levels, such as the 53.0\% conduct rate, depend on where
the instrument set its thresholds.

Fourth, the instrument may know the show. SVU is heavily documented in
the coding model's pretraining data, in fan wikis, recaps, and
commentary on famous episodes, so the coder could in principle recognize
an episode and import outside knowledge, who really did it, how the hour
is remembered, rather than coding the transcript in front of it. Nothing
in the output proves this happened, and nothing rules it out. The
metadata-stripped recode planned for the companion study, which
withholds season and episode identifiers from the coder, doubles as a
test: if the coding shifts when the model can no longer tell which
episode it is reading, contamination was doing work, and the validation
study will treat any such shift as instrument error.

Fifth, the codebook's precision-first design excludes ambiguous cases,
so 521 is best read as a floor. Sixth, the model leaked small amounts of
out-of-vocabulary coding, disclosed in Section 3.4, and the unit of
analysis is the person-episode: eight records are recurring persons
accused in separate episodes, so claims about unique lives should use
the figure of roughly 513. Seventh, the dataset codes no demographics.
Questions about the race, class, and gender of the falsely accused await
the next pass, and the SVU literature makes them pressing (Britto et
al., 2007).

Most of the central proportions survive these limitations comfortably:
the sensitivity band on the episode rate is 2.7 points, the death audit
was adjudicated row by row precisely because it is the most quotable
finding, and the exposure results come from controlled fields with
almost no missing data. The marginal apology rate is the exception, for
the truncation reason above, and I flag it rather than defend it. Its
internal gradient, apology tracking severity and ignoring conduct, is
the part of the finding I would stake something on.

\subsection{7. The Acceptable Casualty}\label{the-acceptable-casualty}

Every long-running procedural teaches its audience an accounting: this
is what an investigation costs, and this is who pays. The account this
study reads out of SVU is specific. The payer is an innocent person in
three of every five hours. The price is usually arrest, a job, a
marriage, or a reputation, and in 32 cases a life. The bill is presented
by the detectives themselves: their inference starts most false
accusations, which a detective show partly guarantees, and, less
forgivably, their disclosure habits accompany the worst damage while
their apologies almost never arrive. And the ledger closes without
acknowledgment 92.1\% of the time.

I set out to measure a narrative convention, and the convention turned
out to have a moral architecture: harm that is shown, so the drama can
use it, and accountability that is withheld, so the heroes can survive
it. Twenty-seven seasons of that architecture, syndicated daily, is a
sustained rehearsal of what wrongful accusation looks like and what
follows it. What follows it, on SVU, is the next case.

The dataset, the codebook, and the audit trail behind this article will
be deposited with the project's OSF repository at submission, and two
extensions are underway: the companion longitudinal study (Orden, in
preparation) and the human validation study the method requires. The
count itself now exists. Anyone who wants to argue that the falsely
accused of SVU are treated better, or worse, than I have described can
check the rows.

\subsection{Acknowledgments}\label{acknowledgments}

Claude (Anthropic) performed the episode coding for this study as the
measurement instrument described in Section 3, applying the fixed
codebook to all 576 transcripts through the Messages Batch API, and
drafted the manuscript text under my direction, which I revised and
verified. I validated every figure in this article against the
underlying dataset, including the adjudicated death audit in Section
4.8, and I take full responsibility for the content. One circularity
should be on the record: the same model family coded the data and
drafted the text, which is a further reason the planned validation
cross-check uses a non-Anthropic model. Episode transcripts were
collected from subslikescript.com.

\subsection{References}\label{references}

Bernabo, L. (2022). Copaganda and post-Floyd TVPD: Broadcast
television's response to policing in 2020. Journal of Communication,
72(4), 488-496.

Britto, S., Hughes, T., Saltzman, K., \& Stroh, C. (2007). Does
``special'' mean young, white and female? Deconstructing the meaning of
``special'' in Law \& Order: Special Victims Unit. Journal of Criminal
Justice and Popular Culture, 14(1), 39-57.

Cole, S. A., \& Dioso-Villa, R. (2009). Investigating the ``CSI effect''
effect: Media and litigation crisis in criminal law. Stanford Law
Review, 61(6), 1335-1373.

Cuklanz, L. M. (2000). Rape on prime time: Television, masculinity, and
sexual violence. University of Pennsylvania Press.

Cuklanz, L. M., \& Moorti, S. (2006). Television's ``new'' feminism:
Prime-time representations of women and victimization. Critical Studies
in Media Communication, 23(4), 302-321.

Dowler, K. (2003). Media consumption and public attitudes toward crime
and justice: The relationship between fear of crime, punitive attitudes,
and perceived police effectiveness. Journal of Criminal Justice and
Popular Culture, 10(2), 109-126.

Eschholz, S., Mallard, M., \& Flynn, S. (2004). Images of prime time
justice: A content analysis of ``NYPD Blue'' and ``Law \& Order.''
Journal of Criminal Justice and Popular Culture, 10(3), 161-180.

Garrett, B. L. (2011). Convicting the innocent: Where criminal
prosecutions go wrong. Harvard University Press.

Gerbner, G., Gross, L., Morgan, M., Signorielli, N., \& Shanahan, J.
(2002). Growing up with television: Cultivation processes. In J. Bryant
\& D. Zillmann (Eds.), Media effects: Advances in theory and research
(2nd ed., pp.~43-67). Lawrence Erlbaum Associates.

Gilardi, F., Alizadeh, M., \& Kubli, M. (2023). ChatGPT outperforms
crowd workers for text-annotation tasks. Proceedings of the National
Academy of Sciences, 120(30), e2305016120.

Goffman, E. (1963). Stigma: Notes on the management of spoiled identity.
Prentice-Hall.

Gross, S. R., Jacoby, K., Matheson, D. J., Montgomery, N., \& Patil, S.
(2005). Exonerations in the United States 1989 through 2003. Journal of
Criminal Law and Criminology, 95(2), 523-560.

Grounds, A. (2004). Psychological consequences of wrongful conviction
and imprisonment. Canadian Journal of Criminology and Criminal Justice,
46(2), 165-182.

Hust, S. J. T., et al.~(2015). Law \& Order, CSI, and NCIS: The
association between exposure to crime drama franchises, rape myth
acceptance, and sexual consent negotiation among college students.
Journal of Health Communication, 20(12), 1369-1381.

Kappeler, V. E., \& Potter, G. W. (2018). The mythology of crime and
criminal justice (5th ed.). Waveland Press.

Karakatsanis, A. (2025). Copaganda: How police and the media manipulate
our news. The New Press.

Kort-Butler, L. A., \& Sittner Hartshorn, K. J. (2011). Watching the
detectives: Crime programming, fear of crime, and attitudes about the
criminal justice system. The Sociological Quarterly, 52(1), 36-55.

Krippendorff, K. (2019). Content analysis: An introduction to its
methodology (4th ed.). Sage.

Lisak, D., Gardinier, L., Nicksa, S. C., \& Cote, A. M. (2010). False
allegations of sexual assault: An analysis of ten years of reported
cases. Violence Against Women, 16(12), 1318-1334.

Loeffler, C. E., Hyatt, J., \& Ridgeway, G. (2019). Measuring
self-reported wrongful convictions among prisoners. Journal of
Quantitative Criminology, 35(2), 259-286.

National Registry of Exonerations. (2026). The National Registry of
Exonerations {[}Database{]}. University of California Irvine; University
of Michigan; Michigan State University. https://exonerationregistry.org

Neuendorf, K. A. (2017). The content analysis guidebook (2nd ed.). Sage.

Orden, T. (in preparation). The retreat of the false accusation plot in
Law \& Order: SVU, 1999-2026. Manuscript.

Podlas, K. (2006). ``The CSI effect'': Exposing the media myth. Fordham
Intellectual Property, Media and Entertainment Law Journal, 16(2),
429-465.

Projansky, S. (2001). Watching rape: Film and television in postfeminist
culture. NYU Press.

Rapping, E. (2003). Law and justice as seen on TV. New York University
Press.

Surette, R. (2015). Media, crime, and criminal justice: Images,
realities, and policies (5th ed.). Cengage Learning.

Törnberg, P. (2023). ChatGPT-4 outperforms experts and crowd workers in
annotating political Twitter messages with zero-shot learning.
arXiv:2304.06588.

Tyler, T. R. (2006). Why people obey the law. Princeton University
Press.

Westervelt, S. D., \& Cook, K. J. (2012). Life after death row:
Exonerees' search for community and identity. Rutgers University Press.

Ziems, C., Held, W., Shaikh, O., Chen, J., Zhang, Z., \& Yang, D.
(2024). Can large language models transform computational social
science? Computational Linguistics, 50(1), 237-291.

\end{document}
